% Clase 23 --- El dibujito del promedio de coseno cuadrado (§1.3).
% El docente lo hace en el pizarrón para justificar <cos^2> = 1/2 sin integrar:
% cos^2 nunca es negativo, oscila al DOBLE de frecuencia y es simetrico respecto
% de la recta 1/2. El area sobre esa recta compensa exactamente la de abajo.
\begin{center}
\begin{tikzpicture}

  \begin{axis}[
    fisfig, width=12cm, height=5.4cm,
    xlabel={$\omega t - \phi$}, ylabel={},
    xmin=0, xmax=6.5, ymin=-1.25, ymax=1.6,
    xtick={1.5708,3.1416,4.7124,6.2832},
    xticklabels={$\frac{\pi}{2}$,$\pi$,$\frac{3\pi}{2}$,$2\pi$},
    ytick={-1,0,0.5,1}, yticklabels={$-1$,$0$,$\tfrac12$,$1$},
    legend style={at={(0.5,-0.45)}, anchor=north, legend columns=3},
  ]
    \addplot[figgray, smooth, samples=200, domain=0:6.5] {cos(deg(x))};
    \addlegendentry{$\cos(\omega t-\phi)$}
    \addplot[figblue, smooth, samples=200, domain=0:6.5]
      {(cos(deg(x)))^2};
    \addlegendentry{$\cos^2(\omega t-\phi)$}
    \addplot[figred, dashed, samples=2, domain=0:6.5] {0.5};
    \addlegendentry{valor medio $=\tfrac12$}

    \draw[figgray, line width=0.6pt,
          {Latex[length=1.6mm,width=1.2mm]}-{Latex[length=1.6mm,width=1.2mm]}]
      (axis cs:3.1416,1.14) -- (axis cs:6.2832,1.14);
    \node[figlblaux, anchor=south, inner sep=1.5pt] at (axis cs:4.7124,1.14)
      {periodo $\pi$: el doble de frecuencia};
  \end{axis}

\end{tikzpicture}\\[0.5em]
{\small\itshape Elevar al cuadrado hace dos cosas a la vez. Primero, borra el
signo: $\cos^2$ nunca baja de cero, y por eso una resistencia disipa energía
en los dos semiciclos, tanto cuando la corriente va como cuando vuelve.
Segundo, duplica la frecuencia, porque los dos máximos del coseno (arriba y
abajo) pasan a ser máximos iguales de $\cos^2$. Lo que hace innecesaria la
integral es la simetría que queda a la vista: la curva es simétrica respecto de
la recta $\tfrac12$, así que lo que sobresale por encima llena exactamente los
huecos de abajo y el promedio en un período es $\tfrac12$ ---el mismo argumento
escrito con fórmulas es $\cos^2x = \tfrac{1+\cos 2x}{2}$, donde el $\cos 2x$
promedia cero. De ahí sale el factor $\tfrac12$ de $\langle P_R\rangle =
\tfrac12 R I_m^2$, y de querer sacárselo de encima nacen los valores eficaces.}
\end{center}
